%-----------------------------------------------%
%                      致谢
%-----------------------------------------------%
\chapter*{致\quad 谢}
\addcontentsline{toc}{chapter}{致~~~~谢}
墨尽笔歇，文成思驻。掩卷回首，求学之路如涉江采芙蓉，风雨兼程中幸得诸君相伴。
此篇拙作不仅凝结了寒窗数载的思考，更承载了难以言尽的感激之情。

致恩师：
师恩如山，润物无声。感谢导师范文博老师，何其有幸，得遇良师，您严谨求真的治学态度，
让我领悟到“治学如琢玉，精益方成器”的真谛；您对学术的赤诚与热忱，必将成为我日后学习与工作的永恒准则。
愿您学术之树常青，桃李满园；盼您身体康泰，万事顺遂。

致挚友：
友谊如星，辉映长夜。
感谢石锐，时光会筛去浮沫，你是被筛出的珍宝，真正的温暖，是知道有人永远在不动声色的陪伴；
感谢张凤，时常用旁观者的清醒点破我的执念，又用“相信你”的底气，给我无尽的鼓励和支持；
感谢徐刘希，“一见如故”的友谊，很幸运遇见“无需讨好、不必伪装”的你，有些人一旦遇见便是永恒的春暖花开；
感谢张坤，初遇时的键盘声是友情的序章，感谢在我生命里敲出温暖的注脚，愿未来我们都能在
热爱里敲出星辰大海；
感谢张振雯和李芙蓉，我青春里的双生星光，做我情绪的“垃圾桶”，做我论文的“进度监工队”，愿我们
永远保持这份“互相麻烦”的底气；
感谢王语嫣、刘超杰、叶少聪、曾爽、姜宇轩和张皓森，“北斗七星”一个都不能少，我们各自闪耀
着不同的光芒，却在同一片夜空下，组成永不迷路的指引；
感谢AtPT Lab 的师兄师姐师弟师妹们：
严啸天师兄实验室的“大家长”坐标，我心中永远的“万能”大师兄；
郑智渊师兄我论文迷雾的引路人，缓解我初遇课题时的不安和焦虑；感谢王福超师兄星巴克味的治愈和
王曙光师兄麻将桌上的“赞助”；
还有师弟师妹们，课题组因你们而温暖明亮，充满活力。
江湖路远，愿我们永远保持"少年击剑更吹箫"的赤诚。 

致亲恩：
亲恩如海，寸草难酬。感谢父母，二十五年养育恩重，从未言说的那些灶台加热的饭菜、车站不舍的目送、微信里的牵挂，
皆是我暗夜里的光。感谢爷爷奶奶姥姥，归家时的期盼、临行前的不舍，用细腻的爱给了我莫大的支持和前行的勇气。
唯愿余生如慈乌绕枝，朝朝暮暮啼暖亲怀，岁岁年年敬报亲恩。

致自己：
这一程山高水长，你终是守住了出发时的那轮明月。愿你永远是那个“笨拙却认真”的少年，永远相信出发时的那轮明月，
终会在岁月的尽头，与你眼中的星河遥遥相认。

书短意长，此情难尽。愿携所得继续前行，以萤火之光，报岁月深恩。唯愿此心似明月，照我前行，亦映照所有曾予我光芒之人。

% \vspace{5mm}

\begin{flushright}
刘月荣\\
二〇二五年五月十八日于西南交通大学
\end{flushright}
